% oss-plantilla: oss-2026.11
% Preámbulo corporativo fijo. No editar en sd-NN.tex: exportar_latex.py verificar lo compara byte a byte.
% Para cambiar el formato, editar esta plantilla u oss.sty y subir la versión de la primera línea.
\PassOptionsToPackage{unicode}{hyperref}
\PassOptionsToPackage{hyphens}{url}
\documentclass[11pt]{article}
\usepackage{amsmath,amssymb}
\usepackage{iftex}
\ifPDFTeX
  % Prism y otros editores en la nube: pdfLaTeX con Helvetica como familia por defecto.
  \usepackage[T1]{fontenc}
  \usepackage[utf8]{inputenc}
  \usepackage{textcomp}
  \usepackage[scaled=0.92]{helvet}
  \renewcommand{\familydefault}{\sfdefault}
  \DeclareUnicodeCharacter{2264}{\ensuremath{\leq}}
  \DeclareUnicodeCharacter{2265}{\ensuremath{\geq}}
  \DeclareUnicodeCharacter{2260}{\ensuremath{\neq}}
  \DeclareUnicodeCharacter{2248}{\ensuremath{\approx}}
  \DeclareUnicodeCharacter{2190}{\ensuremath{\leftarrow}}
  \DeclareUnicodeCharacter{2192}{\ensuremath{\rightarrow}}
  \DeclareUnicodeCharacter{25CF}{\ensuremath{\bullet}}
  \DeclareUnicodeCharacter{2713}{\ensuremath{\checkmark}}
\else
  % XeLaTeX/LuaLaTeX (compilación oficial local): DejaVu Sans.
  \usepackage{fontspec}
  \defaultfontfeatures{Ligatures=TeX,Scale=MatchLowercase}
  \setmainfont{DejaVu Sans}
  \setsansfont{DejaVu Sans}
  \setmonofont{DejaVu Sans Mono}
\fi
\usepackage[letterpaper,margin=20mm]{geometry}
\usepackage{microtype}
\usepackage{parskip}
\usepackage{xcolor}
\usepackage{longtable,booktabs,array}
\usepackage{calc}
\usepackage{etoolbox}
\makeatletter
\patchcmd\longtable{\par}{\if@noskipsec\mbox{}\fi\par}{}{}
\makeatother
\usepackage{footnotehyper}
\makesavenoteenv{longtable}
\usepackage{graphicx}
\makeatletter
\newsavebox\pandoc@box
\newcommand*\pandocbounded[1]{%
  \sbox\pandoc@box{#1}%
  \Gscale@div\@tempa{\textheight}{\dimexpr\ht\pandoc@box+\dp\pandoc@box\relax}%
  \Gscale@div\@tempb{\linewidth}{\wd\pandoc@box}%
  \ifdim\@tempb\p@<\@tempa\p@\let\@tempa\@tempb\fi
  \ifdim\@tempa\p@<\p@\scalebox{\@tempa}{\usebox\pandoc@box}%
  \else\usebox{\pandoc@box}%
  \fi}
\def\maxwidth{\ifdim\Gin@nat@width>\linewidth\linewidth\else\Gin@nat@width\fi}
\def\maxheight{\ifdim\Gin@nat@height>\textheight\textheight\else\Gin@nat@height\fi}
\makeatother
\setkeys{Gin}{width=\maxwidth,height=\maxheight,keepaspectratio}
\makeatletter
\def\fps@figure{htbp}
\makeatother
\usepackage[normalem]{ulem}
\providecommand{\st}[1]{\sout{#1}}
\usepackage{fancyvrb}
\setlength{\emergencystretch}{3em}
\providecommand{\tightlist}{%
  \setlength{\itemsep}{0pt}\setlength{\parskip}{0pt}}
\providecommand{\passthrough}[1]{#1}
\setcounter{secnumdepth}{-\maxdimen}
\ifPDFTeX
  \usepackage[spanish]{babel}
\else
  \usepackage[bidi=default]{babel}
  \babelprovide[main,import]{spanish}
\fi
\let\LanguageShortHands\languageshorthands
\def\languageshorthands#1{}
\usepackage{bookmark}
\usepackage{xurl}
\urlstyle{same}
\hypersetup{pdflang={es-CL},pdfcreator={Only Simple Solutions - exportar_latex.py}}
\usepackage{oss}
\begin{document}

\hypersetup{pageanchor=false}
\ossCover{Anexos}{Subdocumento 6}
\hypersetup{pageanchor=true}
\makeatletter\renewcommand{\@pnumwidth}{2.2em}\makeatother
\tableofcontents
\clearpage

\renewcommand{\tablename}{Tabla}
\section*{Lista de tablas}
\addcontentsline{toc}{section}{Lista de tablas}

\noindent\hangindent=2em\hyperref[tab:a-1-1]{Tabla A.1.1: Planilla del registro de interesados}\dotfill\pageref{tab:a-1-1}\par\vspace{0.4\baselineskip}
\noindent\hangindent=2em\hyperref[tab:a-1-2]{Tabla A.1.2: Instrucciones de llenado del registro de interesados}\dotfill\pageref{tab:a-1-2}\par\vspace{0.4\baselineskip}
\noindent\hangindent=2em\hyperref[tab:a-2-1]{Tabla A.2.1: Planilla de la matriz de evaluación del involucramiento de los interesados}\dotfill\pageref{tab:a-2-1}\par\vspace{0.4\baselineskip}
\noindent\hangindent=2em\hyperref[tab:a-2-2]{Tabla A.2.2: Escala de niveles de involucramiento}\dotfill\pageref{tab:a-2-2}\par\vspace{0.4\baselineskip}
\clearpage

\section{A. Anexos}\label{a-anexos}

Este anexo consolida las planillas de gestión de interesados referidas en la sección 6.1.4 del Subdocumento 6. Ambas se entregan sin completar de forma intencional: se llenarán durante el levantamiento del mes 1, con las entrevistas, encuestas y talleres descritos en la sección 6.1.4.1, y se mantendrán actualizadas durante todo el contrato en el PMIS definido en la sección 6.1.5.1. Los códigos I-01 a I-26 corresponden a los interesados identificados en la Figura 1 del Subdocumento 6.

\subsection{A.1. Planilla del registro de interesados}\label{a-1-planilla-del-registro-de-interesados}

La información obtenida en el levantamiento se documentará en el registro de interesados. Las columnas siguen los tres bloques de información que el PMBOK define para el registro de interesados: identificación, evaluación y clasificación (PMI, 2017). El bloque de identificación permitirá contactar a cada representante; el de evaluación registrará qué espera del proyecto y cuánto puede influir en él; y el de clasificación aplicará los cuatro criterios descritos en la sección 6.1.4.1 del Subdocumento 6. El registro se mantendrá actualizado durante todo el contrato y tendrá acceso restringido en el PMIS, porque las evaluaciones de actitud son información sensible: solo lo consultarán el jefe de proyecto, el líder de implantación y gestión del cambio y el patrocinador de Ancoa.

\clearpage
\subsubsection{A.1.1. Planilla del registro de interesados}\label{a-1-1-planilla-del-registro-de-interesados}

A continuación, se presenta la planilla con la que se levantará el registro de interesados.

\renewcommand{\thetable}{A.1.1}
\begin{longtable}[]{@{}
  >{\raggedright\arraybackslash}p{(\columnwidth - 10\tabcolsep) * \real{0.1000}}
  >{\raggedright\arraybackslash}p{(\columnwidth - 10\tabcolsep) * \real{0.2400}}
  >{\raggedright\arraybackslash}p{(\columnwidth - 10\tabcolsep) * \real{0.2100}}
  >{\raggedright\arraybackslash}p{(\columnwidth - 10\tabcolsep) * \real{0.2300}}
  >{\raggedright\arraybackslash}p{(\columnwidth - 10\tabcolsep) * \real{0.2200}}@{}}
\caption{Planilla del registro de interesados.}\label{tab:a-1-1}\\
\toprule\noalign{}
\multicolumn{2}{@{}>{\raggedright\arraybackslash}p{(\columnwidth - 10\tabcolsep) * \real{0.3400} + 2\tabcolsep}}{\textbf{Nombre del proyecto:}} & \multicolumn{2}{>{\raggedright\arraybackslash}p{(\columnwidth - 10\tabcolsep) * \real{0.4400} + 2\tabcolsep}}{\textbf{Código del proyecto:}} & \multicolumn{1}{>{\raggedright\arraybackslash}p{(\columnwidth - 10\tabcolsep) * \real{0.2200} + 0\tabcolsep}}{\textbf{Fecha de actualización:}} \\[10pt]
\midrule\noalign{}
\begin{minipage}[b]{\linewidth}\raggedright
\textbf{ID}
\end{minipage} & \begin{minipage}[b]{\linewidth}\raggedright
\textbf{Identificación}\\[1pt]{\footnotesize nombre, cargo, correo y teléfono}
\end{minipage} & \begin{minipage}[b]{\linewidth}\raggedright
\textbf{Rol y grupo de interés}\\[1pt]{\footnotesize rol en el proyecto y grupo}
\end{minipage} & \begin{minipage}[b]{\linewidth}\raggedright
\textbf{Evaluación}\\[1pt]{\footnotesize requerimientos, expectativas e influencia}
\end{minipage} & \begin{minipage}[b]{\linewidth}\raggedright
\textbf{Clasificación}\\[1pt]{\footnotesize interno o externo, actor o no operativo, contexto y poder/interés}
\end{minipage} \\
\midrule\noalign{}
\endfirsthead
\toprule\noalign{}
\multicolumn{2}{@{}>{\raggedright\arraybackslash}p{(\columnwidth - 10\tabcolsep) * \real{0.3400} + 2\tabcolsep}}{\textbf{Nombre del proyecto:}} & \multicolumn{2}{>{\raggedright\arraybackslash}p{(\columnwidth - 10\tabcolsep) * \real{0.4400} + 2\tabcolsep}}{\textbf{Código del proyecto:}} & \multicolumn{1}{>{\raggedright\arraybackslash}p{(\columnwidth - 10\tabcolsep) * \real{0.2200} + 0\tabcolsep}}{\textbf{Fecha de actualización:}} \\[10pt]
\midrule\noalign{}
\begin{minipage}[b]{\linewidth}\raggedright
\textbf{ID}
\end{minipage} & \begin{minipage}[b]{\linewidth}\raggedright
\textbf{Identificación}\\[1pt]{\footnotesize nombre, cargo, correo y teléfono}
\end{minipage} & \begin{minipage}[b]{\linewidth}\raggedright
\textbf{Rol y grupo de interés}\\[1pt]{\footnotesize rol en el proyecto y grupo}
\end{minipage} & \begin{minipage}[b]{\linewidth}\raggedright
\textbf{Evaluación}\\[1pt]{\footnotesize requerimientos, expectativas e influencia}
\end{minipage} & \begin{minipage}[b]{\linewidth}\raggedright
\textbf{Clasificación}\\[1pt]{\footnotesize interno o externo, actor o no operativo, contexto y poder/interés}
\end{minipage} \\
\midrule\noalign{}
\endhead
\bottomrule\noalign{}
\endlastfoot
\rule{0pt}{0.85cm} & & & & \\
\midrule\noalign{}
\rule{0pt}{0.85cm} & & & & \\
\midrule\noalign{}
\rule{0pt}{0.85cm} & & & & \\
\midrule\noalign{}
\rule{0pt}{0.85cm} & & & & \\
\midrule\noalign{}
\rule{0pt}{0.85cm} & & & & \\
\midrule\noalign{}
\rule{0pt}{0.85cm} & & & & \\
\midrule\noalign{}
\rule{0pt}{0.85cm} & & & & \\
\midrule\noalign{}
\rule{0pt}{0.85cm} & & & & \\
\end{longtable}

\subsubsection{A.1.2. Instrucciones de llenado del registro de interesados}\label{a-1-2-instrucciones-de-llenado-del-registro-de-interesados}

A continuación, se detalla el contenido que se registrará en cada columna de la planilla.

\renewcommand{\thetable}{A.1.2}
\begin{longtable}[]{@{}
  >{\raggedright\arraybackslash}p{(\columnwidth - 4\tabcolsep) * \real{0.2600}}
  >{\raggedright\arraybackslash}p{(\columnwidth - 4\tabcolsep) * \real{0.7400}}@{}}
\caption{Instrucciones de llenado del registro de interesados.}\label{tab:a-1-2}\\
\toprule\noalign{}
\begin{minipage}[b]{\linewidth}\raggedright
\textbf{Columna}
\end{minipage} & \begin{minipage}[b]{\linewidth}\raggedright
\textbf{Contenido que se registrará}
\end{minipage} \\
\midrule\noalign{}
\endfirsthead
\toprule\noalign{}
\begin{minipage}[b]{\linewidth}\raggedright
\textbf{Columna}
\end{minipage} & \begin{minipage}[b]{\linewidth}\raggedright
\textbf{Contenido que se registrará}
\end{minipage} \\
\midrule\noalign{}
\endhead
\bottomrule\noalign{}
\endlastfoot
\textbf{ID} & Código del interesado, de I-01 a I-26, según la Figura 1 del Subdocumento 6. \\
\addlinespace[2pt]
\textbf{Identificación} & Nombre, cargo, correo y teléfono del representante designado. \\
\addlinespace[2pt]
\textbf{Rol y grupo de interés} & Rol que cumple en el proyecto y grupo de interés al que pertenece. \\
\addlinespace[2pt]
\textbf{Evaluación} & Requerimientos, expectativas e influencia sobre el proyecto. \\
\addlinespace[2pt]
\textbf{Clasificación} & Interno o externo; actor o no operativo; contexto (retail, emisor, frontera retail--emisor o plataforma compartida); y cuadrante de poder e interés. \\
\end{longtable}

\clearpage
\subsection{A.2. Planilla de la matriz de evaluación del involucramiento}\label{a-2-planilla-de-la-matriz-de-evaluacion-del-involucramiento}

La matriz de evaluación del involucramiento compara el nivel actual de compromiso de cada interesado con el nivel que el proyecto necesita para tener éxito (PMI, 2017). La planilla se entrega sin completar de forma intencional: en el mes 1, el nivel actual de cada interesado se evaluará con las entrevistas, encuestas y talleres descritos en la sección 6.1.4.1 del Subdocumento 6, y el nivel deseado se acordará con los representantes de Ancoa antes del hito H1. Los niveles se escribirán con su nombre completo, para no confundir la letra D de desconocedor con la de deseado. La matriz se actualizará al cierre de cada intervalo de planificación.

\subsubsection{A.2.1. Planilla de la matriz de evaluación del involucramiento}\label{a-2-1-planilla-de-la-matriz-de-evaluacion-del-involucramiento}

A continuación, se presenta la planilla de la matriz de evaluación del involucramiento de los interesados.

\renewcommand{\thetable}{A.2.1}
\begin{longtable}[]{@{}
  >{\raggedright\arraybackslash}p{(\columnwidth - 10\tabcolsep) * \real{0.1000}}
  >{\raggedright\arraybackslash}p{(\columnwidth - 10\tabcolsep) * \real{0.2400}}
  >{\raggedright\arraybackslash}p{(\columnwidth - 10\tabcolsep) * \real{0.2100}}
  >{\raggedright\arraybackslash}p{(\columnwidth - 10\tabcolsep) * \real{0.2100}}
  >{\raggedright\arraybackslash}p{(\columnwidth - 10\tabcolsep) * \real{0.2400}}@{}}
\caption{Planilla de la matriz de evaluación del involucramiento de los interesados.}\label{tab:a-2-1}\\
\toprule\noalign{}
\multicolumn{2}{@{}>{\raggedright\arraybackslash}p{(\columnwidth - 10\tabcolsep) * \real{0.3400} + 2\tabcolsep}}{\textbf{Fecha de evaluación:}} & \multicolumn{3}{>{\raggedright\arraybackslash}p{(\columnwidth - 10\tabcolsep) * \real{0.6600} + 4\tabcolsep}}{\textbf{Intervalo de planificación:}} \\[10pt]
\midrule\noalign{}
\begin{minipage}[b]{\linewidth}\raggedright
\textbf{ID}
\end{minipage} & \begin{minipage}[b]{\linewidth}\raggedright
\textbf{Interesado}
\end{minipage} & \begin{minipage}[b]{\linewidth}\raggedright
\textbf{Nivel actual}\\[1pt]{\footnotesize desconocedor, reticente, neutral, de apoyo o líder}
\end{minipage} & \begin{minipage}[b]{\linewidth}\raggedright
\textbf{Nivel deseado}\\[1pt]{\footnotesize nivel que el proyecto necesita e hito en que debe alcanzarlo}
\end{minipage} & \begin{minipage}[b]{\linewidth}\raggedright
\textbf{Acción de involucramiento}\\[1pt]{\footnotesize acción principal para cerrar la brecha}
\end{minipage} \\
\midrule\noalign{}
\endfirsthead
\toprule\noalign{}
\multicolumn{2}{@{}>{\raggedright\arraybackslash}p{(\columnwidth - 10\tabcolsep) * \real{0.3400} + 2\tabcolsep}}{\textbf{Fecha de evaluación:}} & \multicolumn{3}{>{\raggedright\arraybackslash}p{(\columnwidth - 10\tabcolsep) * \real{0.6600} + 4\tabcolsep}}{\textbf{Intervalo de planificación:}} \\[10pt]
\midrule\noalign{}
\begin{minipage}[b]{\linewidth}\raggedright
\textbf{ID}
\end{minipage} & \begin{minipage}[b]{\linewidth}\raggedright
\textbf{Interesado}
\end{minipage} & \begin{minipage}[b]{\linewidth}\raggedright
\textbf{Nivel actual}\\[1pt]{\footnotesize desconocedor, reticente, neutral, de apoyo o líder}
\end{minipage} & \begin{minipage}[b]{\linewidth}\raggedright
\textbf{Nivel deseado}\\[1pt]{\footnotesize nivel que el proyecto necesita e hito en que debe alcanzarlo}
\end{minipage} & \begin{minipage}[b]{\linewidth}\raggedright
\textbf{Acción de involucramiento}\\[1pt]{\footnotesize acción principal para cerrar la brecha}
\end{minipage} \\
\midrule\noalign{}
\endhead
\bottomrule\noalign{}
\endlastfoot
\rule{0pt}{0.85cm} & & & & \\
\midrule\noalign{}
\rule{0pt}{0.85cm} & & & & \\
\midrule\noalign{}
\rule{0pt}{0.85cm} & & & & \\
\midrule\noalign{}
\rule{0pt}{0.85cm} & & & & \\
\midrule\noalign{}
\rule{0pt}{0.85cm} & & & & \\
\midrule\noalign{}
\rule{0pt}{0.85cm} & & & & \\
\midrule\noalign{}
\rule{0pt}{0.85cm} & & & & \\
\midrule\noalign{}
\rule{0pt}{0.85cm} & & & & \\
\end{longtable}

\clearpage
\subsubsection{A.2.2. Escala de niveles de involucramiento}\label{a-2-2-escala-de-niveles-de-involucramiento}

A continuación, se detalla la escala de niveles para completar las columnas de nivel actual y nivel deseado.

\renewcommand{\thetable}{A.2.2}
\begin{longtable}[]{@{}
  >{\raggedright\arraybackslash}p{(\columnwidth - 4\tabcolsep) * \real{0.2000}}
  >{\raggedright\arraybackslash}p{(\columnwidth - 4\tabcolsep) * \real{0.8000}}@{}}
\caption{Escala de niveles de involucramiento.}\label{tab:a-2-2}\\
\toprule\noalign{}
\begin{minipage}[b]{\linewidth}\raggedright
\textbf{Nivel}
\end{minipage} & \begin{minipage}[b]{\linewidth}\raggedright
\textbf{Descripción}
\end{minipage} \\
\midrule\noalign{}
\endfirsthead
\toprule\noalign{}
\begin{minipage}[b]{\linewidth}\raggedright
\textbf{Nivel}
\end{minipage} & \begin{minipage}[b]{\linewidth}\raggedright
\textbf{Descripción}
\end{minipage} \\
\midrule\noalign{}
\endhead
\bottomrule\noalign{}
\endlastfoot
\textbf{Desconocedor} & No conoce el proyecto ni los impactos que puede generar. \\
\addlinespace[2pt]
\textbf{Reticente} & Conoce el proyecto y sus implicancias, pero se resiste al cambio. \\
\addlinespace[2pt]
\textbf{Neutral} & Conoce el proyecto, pero no manifiesta apoyo ni oposición. \\
\addlinespace[2pt]
\textbf{De apoyo} & Comprende el alcance y los impactos del proyecto y colabora con sus actividades y resultados. \\
\addlinespace[2pt]
\textbf{Líder} & Está plenamente informado y comprometido, y promueve activamente el éxito del proyecto dentro de la organización. \\
\end{longtable}

\ossFinalPage

\end{document}
